\section{\revise{Prior Work}}

\revise{Before describing our own
contributions, we show how related works approached
the same problems  
before this dissertation:
security analysis for smart contracts, and 
security of LLM-assisted code
generation. 
For each, we note what earlier work achieved,
why it was not enough to
prevent the incidents of Section~\ref{sec:examples}, and 
how this dissertation
advances the field.}

\subsection{\revise{Securing Smart Contracts}}

\revise{Most prior work aims to find 
vulnerabilities \emph{before} a contract is
deployed. Static-analysis and symbolic-execution tools such as
Oyente~\cite{luu2016making}, Securify~\cite{tsankov2018securify}, and
Slither~\cite{feist2019slither}, along with fuzzers such as
ContractFuzzer~\cite{jiang2018contractfuzzer} and 
ItyFuzz~\cite{shou2023ityfuzz},
detect known vulnerability classes like re-entrancy and 
integer overflow. Yet this
was not enough to prevent incidents like the 
Harvest exploit of
Section~\ref{subsec:example3}: the 
costliest DeFi attacks exploit \emph{design
flaws}, in which every individual operation is 
permitted and only their combination
is harmful, so nothing looks like a bug; 
they span many contracts and long attack vectors,
defeating single-contract analysis; and 
pre-deployment tools cannot help a contract
that is already live. 
A closer line moves protection to runtime, 
learning or
enforcing invariants~\cite{liu2022learning, liu2022invcon, liu2024automated,
deng2024safeguarding} or restricting control 
flow~\cite{rodler2018sereum, li2020securing,
spherex}, but each covers only a narrow set of properties, often needs
manually specified policies, and cannot patch a contract after deployment.}

\revise{This dissertation advances the deployment 
side on exactly these gaps.
Chapter~4 (\textsc{CrossGuard}) targets the cross-contract 
design flaws that defeat
single-contract analysis, automatically adding 
control-flow-integrity checks that
block illegitimate interaction patterns with low 
overhead and no manual
specification. Chapter~5 (\textsc{Trace2Inv}) learns a 
broad range of behavioral
invariants from a contract's own transaction history and 
enforces them as runtime
guards that apply both before and after deployment.}

\subsection{\revise{Securing LLM-Assisted Code Generation}}

\revise{Before this dissertation, the study of insecure AI-generated code largely
fell into three lines: 
inducing harmful code at \emph{inference time} by feeding the
model malicious external content~\cite{zeng2025inducing, yi2025benchmarking};
poisoning the \emph{training pipeline}~\cite{carlini2024poisoning, 
zhao2025data,
jiang2024turning}, 
mostly studied for natural-language rather than code outputs; and
benchmarking how often generated code contains 
classic weaknesses in Common
Weakness Enumeration 
categories~\cite{yang2024seccodeplt, debenedetti2023a}. As
assistants grew into retrieval-augmented agents, 
a parallel line studied poisoning of
retrieved knowledge and early defenses~\cite{zou2025poisonedrag, 
chen2024agentpoison,
zhou2025trustrag, wang2025astute}, but evaluated 
mainly on simple question answering
rather than on coding agents that can execute code. 
None of this established the
threat behind the incident of Section~\ref{subsec:example1}: 
that a \emph{production}
model, with no external access at inference time, 
reproduces a real-world harmful
artifact, such as a live scam endpoint, in response to an 
ordinary developer prompt,
most likely because that artifact was 
passively absorbed into its weights.}

\revise{This dissertation advances 
the development side by closing these gaps.
Chapter~2 moves upstream to coding agents, 
developing defenses against search- and
documentation-level poisoning, where a successful 
attack is most consequential
because the agent can act on the poisoned result. 
Chapter~3 (\textsc{Scam2Prompt})
audits production LLMs with an automated 
framework that issues benign-looking coding
requests at scale and measures how often the 
generated code embeds malicious
endpoints, exposing passive, already-present 
poisoning rather than a new
inference-time attack.}
